% !TeX root = ../../report.tex
\subsection{Cancer Medicine and the Shift to Precision Pathology}\label{sec:cancer-medicine}

Cancer is a disease in which some of the body’s cells grow uncontrollably and spread to other parts of the body, and there are more than 100 types. Many cancers form solid tumours, but blood cancers such as leukaemias generally do not. Cancer is a genetic disease caused by changes to genes, which can arise from errors during cell division, damage from environmental exposures such as tobacco smoke and ultraviolet radiation, or inheritance \citep{nci2026}. Treatment was traditionally based on the anatomical site of the cancer. Precision medicine instead aims to tailor treatment to the individual, and precision pathology contributes by combining histological and molecular profiling of the patient’s cancer \citep{klauschen2024}. This requires integrating high-dimensional molecular data, such as gene-expression profiles and mutation status, to predict how a specific cancer will respond to therapy \citep{klauschen2024}. Machine learning models are now widely used for tasks such as predicting patient outcomes and drug sensitivity \citep{baptista2021}.

\subsubsection{Genomic Resources, Cell-Line Data and the Trust Gap}\label{sec:cell-line-trust-gap}

Progress in this area depends on large genomic resources such as the Genomics of Drug Sensitivity in Cancer (GDSC), which screens drugs against hundreds of cancer cell lines and records molecular profiles for each \citep{yang2013,baptista2021}. Testing every candidate drug on patients is impractical and costly, so cancer cell lines are used as proxies. These are cells grown in the laboratory that retain much of the genetic and molecular character of the tumours they came from \citep{yang2013}. Drug response is summarised by the half-maximal inhibitory concentration (IC50), the drug concentration that reduces the viability of a cell population by 50\%, so a lower IC50 means a more sensitive cell line \citep{yang2013}. Cell lines are an imperfect proxy, because translating predictions from cell lines to patients remains difficult \citep{wu2025}. The findings of this report therefore concern cell lines only.

Deep neural networks predict drug response well from such data but largely operate as black boxes \citep{partin2023,klauschen2024}. Models trained on high-dimensional molecular data can easily find correlations that are spurious and do not generalise \citep{klauschen2024}, so a clinician cannot easily tell whether a prediction rests on genuine oncogenic biology or on such artefacts. Klauschen et al.\ \citeyearpar{klauschen2024} argue that explainable AI is needed to make the black-box character of conventional AI transparent, to detect flawed prediction strategies, and to support regulatory approval and scientific insight.

Unreliable explanations can also originate in the data. Batch effects are non-biological technical variations introduced by different laboratories, instruments, assessment pipelines, or experimental conditions over time \citep{yu2024,goh2022}, and when data come from several sources such differences are likely to bias downstream analyses \citep{partin2023}. A model can exploit systematic differences of this kind as shortcuts. In pathology images, for example, it may recognise the colour signature of each hospital’s slide scanner and so predict correctly for the wrong reason, which is known as Clever Hans behaviour \citep{klauschen2024}. Such a model is at risk of failing on new data and is of little use for extracting scientific knowledge \citep{klauschen2024}. A rule that is mathematically sound but biologically incoherent, for example one linking a housekeeping gene to tumour suppression, could therefore mislead. This motivates auditing explanations against biology instead of accepting them because they look plausible. Section 6 shows why: a training metric did not predict explanation quality, so the rules had to be assessed directly.